\section{Rank-One Agreement Approximation: What Works and Why (with Detailed Proofs)}
\label{sec:rank_one_agreement}

This section proves two statements that are both critical for the proposed approximation method:

\begin{enumerate}
    \item \textbf{Agreement capture.} The averaging-based rank-one estimator
    $\widehat{G}_1$ recovers the dominant grammar-driven \emph{agreement spike}
    $A_{\mathrm{agr}}$ up to a controlled error.
    \item \textbf{Semantic non-capture (yet learnability).} The same averaging probe
    does \emph{not} extract intra-sequence semantic structure: semantic contributions can be
    large at the token level and remain available to SGD, while still canceling in the
    specific aggregated statistics used by the probe.
\end{enumerate}

Throughout, we make the dependence on the foundational assumptions in
Section~\ref{sec:foundational_assumptions} explicit.

\subsection{Setup: agreement component plus semantic residual}
\label{sec:rank_one_setup}

Let $G\in\mathbb{R}^{d\times d}$ be the gradient matrix of a transformer linear layer.
We decompose
\begin{equation}
\label{eq:rank_one_decomp}
G \;=\; A_{\mathrm{agr}} \;+\; R_{\mathrm{sem}},
\end{equation}
where $A_{\mathrm{agr}}$ is the grammar-driven \emph{agreement component} and $R_{\mathrm{sem}}$
is the semantic/context-specific residual.

\paragraph{Rank-one agreement model.}
Assume the agreement component is rank one:
\begin{equation}
\label{eq:rank_one_A}
A_{\mathrm{agr}} \;=\; \sigma a b^\top,
\qquad
a,b\in\mathbb{R}^d\setminus\{0\},
\qquad
\sigma>0.
\end{equation}

\paragraph{Visibility under the averaging probe.}
Let $\mathbf{1}\in\mathbb{R}^d$ be the all-ones vector and define
\begin{equation}
\label{eq:alpha_beta}
\alpha \triangleq \mathbf{1}^\top a,
\qquad
\beta \triangleq \mathbf{1}^\top b.
\end{equation}
Assume $\alpha\neq 0$ and $\beta\neq 0$, so that $A_{\mathrm{agr}}$ is visible to this probe.

\subsection{Semantic cancellation under aggregation: formal sequence-level proof}
\label{sec:semantic_cancellation_proof}

We now prove a concrete, sequence-level cancellation statement for the probe statistics
$R_{\mathrm{sem}}\mathbf{1}$, $R_{\mathrm{sem}}^\top\mathbf{1}$, and $\mathbf{1}^\top R_{\mathrm{sem}}\mathbf{1}$.

\subsubsection{Batch model and assumptions}

Consider a batch consisting of $S$ sequences. For sequence $s\in\{1,\dots,S\}$,
let $\mathcal{T}_s$ be the set of token indices included from that sequence
(e.g.\ all tokens, or a subset of sufficiently distant positions in long-context settings).
Write the batch gradient as
\begin{equation}
\label{eq:G_seq_sum}
G \;=\; \sum_{s=1}^S \sum_{t\in\mathcal{T}_s} \delta^{(s,t)} \bigl(x^{(s,t)}\bigr)^\top.
\end{equation}
Conceptually decompose each token-level contribution as
\begin{equation}
\label{eq:token_decomp}
\delta^{(s,t)} \bigl(x^{(s,t)}\bigr)^\top
\;=\;
A_{\mathrm{agr}}^{(s,t)} \;+\; R_{\mathrm{sem}}^{(s,t)}.
\end{equation}
Define the sequence-level semantic residual sum
\begin{equation}
\label{eq:Zs_def}
Z_s \;\triangleq\; \sum_{t\in\mathcal{T}_s} R_{\mathrm{sem}}^{(s,t)},
\qquad\text{so that}\qquad
R_{\mathrm{sem}} \;=\; \sum_{s=1}^S Z_s.
\end{equation}

We now formalize the statistical content of Section~\ref{sec:foundational_assumptions}:

\begin{assumption}[Sequence-level semantic independence and zero-mean]
\label{assump:Zs}
Conditioned on model parameters and the batch construction rule, the random matrices
$\{Z_s\}_{s=1}^S$ satisfy:
\begin{enumerate}
\item (\emph{Inter-sequence independence}) $Z_1,\dots,Z_S$ are independent.
\item (\emph{Zero-mean across sequences}) $\mathbb{E}[Z_s]=0$ for each $s$.
\end{enumerate}
For long-context training, the same assumption applies when $\mathcal{T}_s$ is chosen to
include only sufficiently distant positions so that semantic dependence within $\mathcal{T}_s$
is weak, i.e.\ $Z_s$ aggregates approximately independent semantic blocks.
\end{assumption}

\begin{assumption}[Finite second moments of probe responses]
\label{assump:moments}
There exist constants $\sigma_u^2,\sigma_v^2,\sigma_s^2<\infty$ such that for all $s$,
\begin{equation}
\label{eq:moment_bounds}
\mathbb{E}\|Z_s\mathbf{1}\|_2^2 \le \sigma_u^2,
\qquad
\mathbb{E}\|Z_s^\top\mathbf{1}\|_2^2 \le \sigma_v^2,
\qquad
\mathbb{E}\bigl(\mathbf{1}^\top Z_s \mathbf{1}\bigr)^2 \le \sigma_s^2.
\end{equation}
\end{assumption}

\subsubsection{Lemma: cancellation of semantic residual in the probe statistics}

\begin{lemma}[Semantic cancellation under aggregation (detailed)]
\label{lem:semantic_cancellation}
Let $R_{\mathrm{sem}}=\sum_{s=1}^S Z_s$ with $\{Z_s\}$ satisfying
Assumptions~\ref{assump:Zs}--\ref{assump:moments}. Then:
\begin{align}
\label{eq:cancel_mean}
\mathbb{E}[R_{\mathrm{sem}}\mathbf{1}] &= 0,
\qquad
\mathbb{E}[R_{\mathrm{sem}}^\top\mathbf{1}] = 0,
\qquad
\mathbb{E}[\mathbf{1}^\top R_{\mathrm{sem}}\mathbf{1}] = 0,\\
\label{eq:cancel_second_moment}
\mathbb{E}\|R_{\mathrm{sem}}\mathbf{1}\|_2^2 &\le S\,\sigma_u^2,
\qquad
\mathbb{E}\|R_{\mathrm{sem}}^\top\mathbf{1}\|_2^2 \le S\,\sigma_v^2,
\qquad
\mathbb{E}\bigl(\mathbf{1}^\top R_{\mathrm{sem}}\mathbf{1}\bigr)^2 \le S\,\sigma_s^2.
\end{align}
Consequently, for any $\delta\in(0,1)$, with probability at least $1-\delta$,
\begin{equation}
\label{eq:cancel_hp}
\|R_{\mathrm{sem}}\mathbf{1}\|_2 \le \sqrt{\frac{S}{\delta}}\,\sigma_u,\quad
\|R_{\mathrm{sem}}^\top\mathbf{1}\|_2 \le \sqrt{\frac{S}{\delta}}\,\sigma_v,\quad
|\mathbf{1}^\top R_{\mathrm{sem}}\mathbf{1}| \le \sqrt{\frac{S}{\delta}}\,\sigma_s.
\end{equation}
\end{lemma}

\begin{proof}
We prove each claim explicitly.

\paragraph{Mean identities.}
Using \eqref{eq:Zs_def} and linearity of expectation,
\[
\mathbb{E}[R_{\mathrm{sem}}\mathbf{1}]
=
\mathbb{E}\Big[\sum_{s=1}^S Z_s\mathbf{1}\Big]
=
\sum_{s=1}^S \mathbb{E}[Z_s]\mathbf{1}
=
0,
\]
where the last equality uses $\mathbb{E}[Z_s]=0$.
The proofs of $\mathbb{E}[R_{\mathrm{sem}}^\top\mathbf{1}]=0$ and
$\mathbb{E}[\mathbf{1}^\top R_{\mathrm{sem}}\mathbf{1}]=0$ are identical.

\paragraph{Second-moment bound for $R_{\mathrm{sem}}\mathbf{1}$.}
Let $U_s \triangleq Z_s\mathbf{1}\in\mathbb{R}^d$. Then $R_{\mathrm{sem}}\mathbf{1}=\sum_{s=1}^S U_s$.
Expand the squared norm:
\[
\Big\|\sum_{s=1}^S U_s\Big\|_2^2
=
\sum_{s=1}^S \|U_s\|_2^2
+
2\sum_{1\le s<t\le S} U_s^\top U_t.
\]
Taking expectations gives
\[
\mathbb{E}\Big\|\sum_{s=1}^S U_s\Big\|_2^2
=
\sum_{s=1}^S \mathbb{E}\|U_s\|_2^2
+
2\sum_{s<t} \mathbb{E}[U_s^\top U_t].
\]
By independence of $Z_s$ and $Z_t$, $U_s$ and $U_t$ are independent; since $\mathbb{E}[U_s]=0$,
\[
\mathbb{E}[U_s^\top U_t]
=
\mathbb{E}[U_s]^\top \mathbb{E}[U_t]
=
0.
\]
Thus,
\[
\mathbb{E}\|R_{\mathrm{sem}}\mathbf{1}\|_2^2
=
\sum_{s=1}^S \mathbb{E}\|U_s\|_2^2
=
\sum_{s=1}^S \mathbb{E}\|Z_s\mathbf{1}\|_2^2
\le
S\,\sigma_u^2,
\]
using \eqref{eq:moment_bounds}. This proves the first inequality in \eqref{eq:cancel_second_moment}.
The other two second-moment inequalities follow by the same argument with
$V_s \triangleq Z_s^\top\mathbf{1}$ and $W_s \triangleq \mathbf{1}^\top Z_s\mathbf{1}$.

\paragraph{High-probability bounds.}
Apply Markov's inequality to nonnegative random variables:
\[
\mathbb{P}\!\left(\|R_{\mathrm{sem}}\mathbf{1}\|_2 \ge \tau\right)
=
\mathbb{P}\!\left(\|R_{\mathrm{sem}}\mathbf{1}\|_2^2 \ge \tau^2\right)
\le
\frac{\mathbb{E}\|R_{\mathrm{sem}}\mathbf{1}\|_2^2}{\tau^2}
\le
\frac{S\sigma_u^2}{\tau^2}.
\]
Setting $\tau=\sqrt{S\sigma_u^2/\delta}$ yields
$\mathbb{P}(\|R_{\mathrm{sem}}\mathbf{1}\|_2 \ge \sqrt{S/\delta}\,\sigma_u)\le\delta$.
The other two bounds follow identically.
\end{proof}

\subsubsection{From cancellation to the $\varepsilon$-style inequalities}

To relate Lemma~\ref{lem:semantic_cancellation} to the relative inequalities used by the estimator,
we also need that the agreement component reinforces \emph{linearly} with the number of independent
sequences, reflecting shared grammar structure.

\begin{assumption}[Agreement reinforcement in probe marginals]
\label{assump:agreement_reinforce}
There exist constants $m_u,m_v,m_s>0$ such that
\begin{equation}
\label{eq:agreement_linear_scaling}
\|A_{\mathrm{agr}}\mathbf{1}\|_2 \ge S\,m_u,\qquad
\|A_{\mathrm{agr}}^\top\mathbf{1}\|_2 \ge S\,m_v,\qquad
|\mathbf{1}^\top A_{\mathrm{agr}}\mathbf{1}| \ge S\,m_s.
\end{equation}
\end{assumption}

\begin{corollary}[Relative semantic cancellation for the probe statistics]
\label{cor:eps_condition}
Under Lemma~\ref{lem:semantic_cancellation} and Assumption~\ref{assump:agreement_reinforce},
for any $\delta\in(0,1)$, with probability at least $1-\delta$,
\begin{equation}
\label{eq:residual_cancel}
\|R_{\mathrm{sem}}\mathbf{1}\|
\le
\varepsilon_u \|A_{\mathrm{agr}}\mathbf{1}\|,
\qquad
\|R_{\mathrm{sem}}^\top\mathbf{1}\|
\le
\varepsilon_v \|A_{\mathrm{agr}}^\top\mathbf{1}\|,
\qquad
|\mathbf{1}^\top R_{\mathrm{sem}}\mathbf{1}|
\le
\varepsilon_s |\mathbf{1}^\top A_{\mathrm{agr}}\mathbf{1}|,
\end{equation}
where
\[
\varepsilon_u \triangleq \frac{1}{\sqrt{S}}\cdot \frac{\sigma_u}{m_u}\cdot\frac{1}{\sqrt{\delta}},
\qquad
\varepsilon_v \triangleq \frac{1}{\sqrt{S}}\cdot \frac{\sigma_v}{m_v}\cdot\frac{1}{\sqrt{\delta}},
\qquad
\varepsilon_s \triangleq \frac{1}{\sqrt{S}}\cdot \frac{\sigma_s}{m_s}\cdot\frac{1}{\sqrt{\delta}}.
\]
In particular, for fixed $\delta$, each $\varepsilon_\bullet = O(S^{-1/2})$.
\end{corollary}

\begin{proof}
From Lemma~\ref{lem:semantic_cancellation}, with probability at least $1-\delta$,
\[
\|R_{\mathrm{sem}}\mathbf{1}\|_2 \le \sqrt{\frac{S}{\delta}}\,\sigma_u.
\]
From Assumption~\ref{assump:agreement_reinforce},
\[
\|A_{\mathrm{agr}}\mathbf{1}\|_2 \ge S\,m_u.
\]
Divide the first inequality by the second to get
\[
\frac{\|R_{\mathrm{sem}}\mathbf{1}\|_2}{\|A_{\mathrm{agr}}\mathbf{1}\|_2}
\le
\frac{\sqrt{S/\delta}\,\sigma_u}{S\,m_u}
=
\frac{1}{\sqrt{S}}\cdot\frac{\sigma_u}{m_u}\cdot\frac{1}{\sqrt{\delta}}
=
\varepsilon_u.
\]
The other two inequalities are identical.
\end{proof}

\subsection{Agreement spike recovery via the averaging estimator (full proof)}
\label{sec:agreement_recovery}

Define
\begin{equation}
\label{eq:uvs_def}
u \triangleq G\mathbf{1},
\qquad
v \triangleq G^\top \mathbf{1},
\qquad
s \triangleq \mathbf{1}^\top G \mathbf{1},
\end{equation}
and the rank-one estimator
\begin{equation}
\label{eq:Ghat1_def}
\widehat{G}_1 \triangleq \frac{u v^\top}{s}.
\end{equation}

\begin{theorem}[Recovery of the rank-one agreement component]
\label{thm:rank_one_agreement}
Assume \eqref{eq:rank_one_decomp}--\eqref{eq:rank_one_A} and $\alpha\neq 0$, $\beta\neq 0$ in
\eqref{eq:alpha_beta}. Suppose \eqref{eq:residual_cancel} holds for some $0\le\varepsilon<1$
(where for simplicity we take $\varepsilon\triangleq\max\{\varepsilon_u,\varepsilon_v,\varepsilon_s\}$).
Define
\begin{equation}
\label{eq:c_def}
c \triangleq \frac{\sigma^2 \alpha \beta}{s}.
\end{equation}
Then
\begin{align}
\label{eq:rank_one_bound}
\|\widehat{G}_1 - c\,ab^\top\|_2
&\le
\frac{2\varepsilon+\varepsilon^2}{1-\varepsilon}\,\|A_{\mathrm{agr}}\|_2,
\\[3pt]
\label{eq:scale_bound}
\left|\frac{c}{\sigma}-1\right|
&\le
\frac{\varepsilon}{1-\varepsilon}.
\end{align}
\end{theorem}

\begin{proof}
Define perturbations
\[
r_u \triangleq R_{\mathrm{sem}}\mathbf{1},\qquad
r_v \triangleq R_{\mathrm{sem}}^\top\mathbf{1},\qquad
r_s \triangleq \mathbf{1}^\top R_{\mathrm{sem}}\mathbf{1}.
\]
Then using $G=A_{\mathrm{agr}}+R_{\mathrm{sem}}$ and $A_{\mathrm{agr}}=\sigma ab^\top$,
\begin{align*}
u &= G\mathbf{1} = A_{\mathrm{agr}}\mathbf{1} + r_u
= \sigma a(b^\top\mathbf{1}) + r_u
= \sigma\beta a + r_u,\\
v &= G^\top\mathbf{1} = A_{\mathrm{agr}}^\top\mathbf{1} + r_v
= \sigma b(a^\top\mathbf{1}) + r_v
= \sigma\alpha b + r_v,\\
s &= \mathbf{1}^\top G\mathbf{1}
= \mathbf{1}^\top A_{\mathrm{agr}}\mathbf{1} + r_s
= \sigma\alpha\beta + r_s.
\end{align*}

From \eqref{eq:residual_cancel} and $A_{\mathrm{agr}}\mathbf{1}=\sigma\beta a$,
$A_{\mathrm{agr}}^\top\mathbf{1}=\sigma\alpha b$, $\mathbf{1}^\top A_{\mathrm{agr}}\mathbf{1}=\sigma\alpha\beta$,
we have
\[
\|r_u\|\le \varepsilon \|\sigma\beta a\|,
\qquad
\|r_v\|\le \varepsilon \|\sigma\alpha b\|,
\qquad
|r_s|\le \varepsilon |\sigma\alpha\beta|.
\]
In particular,
\begin{equation}
\label{eq:s_lower}
|s|
=
|\sigma\alpha\beta + r_s|
\ge
|\sigma\alpha\beta| - |r_s|
\ge
(1-\varepsilon)|\sigma\alpha\beta|.
\end{equation}

Now expand $\widehat{G}_1$:
\[
\widehat{G}_1
=
\frac{(\sigma\beta a + r_u)(\sigma\alpha b + r_v)^\top}{s}.
\]
Subtract $c\,ab^\top$ with $c=\sigma^2\alpha\beta/s$ and expand:
\[
\widehat{G}_1 - c\,ab^\top
=
\frac{\sigma\beta\, a r_v^\top}{s}
+
\frac{\sigma\alpha\, r_u b^\top}{s}
+
\frac{r_u r_v^\top}{s}.
\]
Using $\|xy^\top\|_2=\|x\|_2\|y\|_2$,
\begin{align*}
\|\widehat{G}_1 - c\,ab^\top\|_2
&\le
\frac{\sigma|\beta|\,\|a\|\,\|r_v\|}{|s|}
+
\frac{\sigma|\alpha|\,\|b\|\,\|r_u\|}{|s|}
+
\frac{\|r_u\|\,\|r_v\|}{|s|}.
\end{align*}
Bound each term using \eqref{eq:s_lower}:
\[
\frac{\sigma|\beta|\,\|a\|\,\|r_v\|}{|s|}
\le
\frac{\sigma|\beta|\,\|a\|}{(1-\varepsilon)|\sigma\alpha\beta|}
\cdot
\varepsilon\|\sigma\alpha b\|
=
\frac{\varepsilon}{1-\varepsilon}\,\sigma\|a\|\|b\|.
\]
Similarly,
\[
\frac{\sigma|\alpha|\,\|b\|\,\|r_u\|}{|s|}
\le
\frac{\varepsilon}{1-\varepsilon}\,\sigma\|a\|\|b\|.
\]
For the third term,
\[
\frac{\|r_u\|\,\|r_v\|}{|s|}
\le
\frac{\varepsilon\|\sigma\beta a\|\cdot \varepsilon\|\sigma\alpha b\|}{(1-\varepsilon)|\sigma\alpha\beta|}
=
\frac{\varepsilon^2}{1-\varepsilon}\,\sigma\|a\|\|b\|.
\]
Summing yields
\[
\|\widehat{G}_1 - c\,ab^\top\|_2
\le
\frac{2\varepsilon+\varepsilon^2}{1-\varepsilon}\,\sigma\|a\|\|b\|.
\]
Since $\|A_{\mathrm{agr}}\|_2=\|\sigma ab^\top\|_2=\sigma\|a\|\|b\|$, this proves \eqref{eq:rank_one_bound}.

Finally,
\[
\frac{c}{\sigma}
=
\frac{\sigma\alpha\beta}{\sigma\alpha\beta+r_s}
=
\frac{1}{1+r_s/(\sigma\alpha\beta)}.
\]
Let $x=r_s/(\sigma\alpha\beta)$. Then $|x|\le\varepsilon$ and
\[
\left|\frac{c}{\sigma}-1\right|
=
\left|\frac{1}{1+x}-1\right|
=
\left|\frac{-x}{1+x}\right|
\le
\frac{|x|}{1-|x|}
\le
\frac{\varepsilon}{1-\varepsilon},
\]
proving \eqref{eq:scale_bound}.
\end{proof}

\subsection{Why the probe does not capture intra-sequence semantics (formal statement)}
\label{sec:semantic_non_capture}

The previous results show that the probe statistics $(G\mathbf{1},G^\top\mathbf{1},\mathbf{1}^\top G\mathbf{1})$
are dominated by agreement. We now formalize the complementary statement: \emph{semantic structure can
remain large and correlated within each sequence, yet still be invisible to the probe.}
This is exactly the sense in which intra-sequence semantics remain learnable.

\paragraph{A precise distinction: first-order probe vs.\ semantic energy.}
The probe only depends on three \emph{first-order} aggregated quantities:
\[
R_{\mathrm{sem}}\mathbf{1},\quad R_{\mathrm{sem}}^\top\mathbf{1},\quad \mathbf{1}^\top R_{\mathrm{sem}}\mathbf{1}.
\]
In contrast, intra-sequence semantic learnability is governed by the \emph{token-level} residuals
$\{R_{\mathrm{sem}}^{(s,t)}\}$ entering SGD updates, and can be quantified by \emph{second-order}
(or energy) measures, e.g.
\[
E_{\mathrm{sem}}
\triangleq
\sum_{s=1}^S \sum_{t\in\mathcal{T}_s} \bigl\|R_{\mathrm{sem}}^{(s,t)}\bigr\|_F^2,
\qquad
\text{or}
\qquad
\sum_{t\in\mathcal{T}_s} R_{\mathrm{sem}}^{(s,t)} \bigl(R_{\mathrm{sem}}^{(s,t)}\bigr)^\top.
\]
The key point is: \emph{first-order cancellation does not imply small energy}.

\begin{lemma}[Cancellation does not imply vanishing semantic energy]
\label{lem:energy_vs_sum}
Let $Z_1,\dots,Z_S$ be independent random matrices with $\mathbb{E}[Z_s]=0$.
Then:
\begin{equation}
\label{eq:sum_small_energy_large}
\mathbb{E}\Big\|\sum_{s=1}^S Z_s\Big\|_F^2
=
\sum_{s=1}^S \mathbb{E}\|Z_s\|_F^2.
\end{equation}
In particular, even though $\sum_s Z_s$ concentrates around $0$ in many directions,
the total semantic energy $\sum_s \mathbb{E}\|Z_s\|_F^2$ can scale linearly in $S$.
\end{lemma}

\begin{proof}
Expand the Frobenius norm:
\[
\Big\|\sum_{s=1}^S Z_s\Big\|_F^2
=
\left\langle \sum_{s=1}^S Z_s,\ \sum_{t=1}^S Z_t\right\rangle_F
=
\sum_{s=1}^S \|Z_s\|_F^2
+
2\sum_{s<t} \langle Z_s, Z_t\rangle_F.
\]
Take expectations. By independence and $\mathbb{E}[Z_s]=0$,
\[
\mathbb{E}\langle Z_s, Z_t\rangle_F
=
\langle \mathbb{E}[Z_s],\mathbb{E}[Z_t]\rangle_F
=
0
\qquad (s\neq t).
\]
Thus
\[
\mathbb{E}\Big\|\sum_{s=1}^S Z_s\Big\|_F^2
=
\sum_{s=1}^S \mathbb{E}\|Z_s\|_F^2,
\]
proving \eqref{eq:sum_small_energy_large}.
\end{proof}

\paragraph{Consequences for intra-sequence semantics and learnability.}
Lemma~\ref{lem:semantic_cancellation} shows that the probe responses
$R_{\mathrm{sem}}\mathbf{1}$, $R_{\mathrm{sem}}^\top\mathbf{1}$, and $\mathbf{1}^\top R_{\mathrm{sem}}\mathbf{1}$
scale like $O(\sqrt{S})$ (and hence become negligible relative to agreement, which is $\Theta(S)$).
Lemma~\ref{lem:energy_vs_sum} shows that simultaneously the semantic energy contained in the
sequence-level residuals $\{Z_s\}$ can scale as $\Theta(S)$ and need not vanish.

Therefore, semantic structure can remain large and correlated within each sequence (and thus
useful for learning), while still being \emph{invisible to the specific averaging probe} used to
estimate the agreement spike. This is exactly the desired behavior: the probe extracts the
grammar-driven spike while leaving intra-sequence semantic correlations untouched for SGD to learn.

Although the analysis in this section is framed in terms of the gradient matrix G, the core mechanism—removal of a dominant shared component while preserving high-dimensional variability—is representation-agnostic. Appendix A formalizes this mechanism in a canonical high-dimensional model, where each column corresponds to a token and decomposes into a shared component plus token-specific residuals. The appendix proves that removing the shared mean direction preserves almost all variance of the residuals, thereby justifying the semantic-preservation claim made here for gradients.